\section{Experiments}
\label{sec:experiments}

Our experiments serve two goals. We first establish that NESS is an effective model for learning under missing node features: it outperforms baselines spanning every category of the problem (Section~\ref{sec:exp-main}), and its advantage widens as features become scarcer (Section~\ref{sec:exp-rate}). We then turn to the study that motivates the paper, asking not merely whether self-supervision helps but \emph{when} and \emph{why}. We show that its benefit is conditional: only some objectives help (Section~\ref{sec:exp-objectives}), for reasons visible in their mechanism and training dynamics (Section~\ref{sec:exp-dynamics}); the benefit emerges primarily under severe missingness (Section~\ref{sec:exp-missingness-regime}) and on dense rather than sparse binary features (Section~\ref{sec:exp-regime}). Finally, we justify NESS's design choices and establish its transferability and scalability (Sections~\ref{sec:exp-design}--\ref{sec:exp-scalability}). Throughout, all learnable methods are compared under the identical protocol described in Section~\ref{sec:setup}.
